\documentclass[11pt,a4paper]{article}
\usepackage[margin=1.1in]{geometry}
\usepackage{amsmath,amssymb,amsthm}
\usepackage{xcolor}
\usepackage{booktabs}
\usepackage{enumitem}
\usepackage[colorlinks=true,linkcolor=blue!50!black,urlcolor=blue!50!black]{hyperref}
\usepackage{titlesec}
\titleformat{\section}{\large\bfseries}{\thesection}{0.8em}{}
\titleformat{\subsection}{\normalsize\bfseries}{\thesubsection}{0.7em}{}

\newcommand{\lift}[1]{\noindent\textbf{Lift:} #1.}
\newcommand{\gives}{\paragraph{What it gives.}}
\newcommand{\form}{\paragraph{Formalization.}}
\newcommand{\proofpath}{\paragraph{Proof path.}}

\title{\textbf{The Provable-Results Portfolio}\\[4pt]
\large Eighteen theorems, bounds, and experiments I can carry end-to-end,\\ with lift estimates and working drafts\\[2pt]
\normalsize For \emph{The Harbor, the Person, and the Economy} (Owens, v1.0)}
\author{Prepared by Claude for Erich Owens}
\date{August 15, 2026}

\begin{document}
\maketitle

\noindent\textbf{Scope and honesty rule.} This document enumerates every result from our exchanges that I judge myself capable of developing \emph{rigorously} --- meaning I can state the model, carry the proof or derivation myself, and tell you exactly which assumptions bear the load --- as opposed to results I could only sketch or survey. Each entry gives the lift (my estimate of focused working time to a polished, checkable draft), a formalization, the actual proof path with its crux, and the payoff to the volume. A closing section lists what I \emph{decline} to claim, in the volume's own spirit. Tier A items are days each; Tier B items are roughly a week each; Tier C items are multi-week. Several Tier A items are independent enough to batch.

\tableofcontents

\section{Tier A --- days each}

\subsection{A1. The Split-Digest Theorem (comonotonicity characterization)}
\lift{1--2 days; the proof is short, the care is in the definitions}

\form Define a \emph{read-surface} as a triple $(X, \Sigma, \pi)$: substrate items $X$, per-item evidence $\Sigma$, and a projection $\pi$ producing a bounded rendering. A \emph{reader} $r$ is a loss function $L_r(\text{rendered set}, \text{true states})$; a scalar head $h: X \to \mathbb{R}$ \emph{serves} $r$ if rank-ordering by $h$ realizes $r$'s optimal rendering at every budget. Two concrete readers: the \emph{successor} (loss $=$ task-continuation value forgone; fit-shaped) and the \emph{operator} (loss $=$ regret-if-ignored; risk-shaped). The theorem: a single head serves both readers \emph{iff} their induced preference orders are comonotone on the item domain --- i.e., agree on every pair. The negative direction is witnessed by two constructible items: a routine-but-essential working state (high continuation value, negligible regret --- the successor must reload it; the operator need never see it) and an anomalously abandoned dead end with irreversible side effects (zero continuation value --- the successor should \emph{not} reload it --- maximal regret). Non-comonotone; hence no shared head; hence Thm.~I.7.1 is a special case of a general principle over every read-surface: digests, checkpoints, resurrection handoffs, briefings.

\proofpath The order-reversal argument is your Thm.~I.7.1's, generalized: strict monotone transforms preserve pairwise order, so a shared head forces comonotonicity; conversely, comonotone orders admit a common refinement, which any strictly increasing reparameterization serves --- giving the \emph{iff}, which the original theorem lacks. The quantitative addendum is where new content lives: apply Thm.~I.6.1 per reader. Each reader has its own critical set $S_r$; a \emph{joint} zero-miss digest must satisfy the floor for $S_{\mathrm{op}} \cup S_{\mathrm{succ}}$, and when the sets are nearly disjoint the joint floor approaches the \emph{sum} of the individual floors --- a single stored compaction pays both bills at once or silently defaults on one. Corollary: compaction must be computed per reader \emph{from durable artifacts}; storing only one compaction is irreversible information loss for the other reader, which is exactly the ``never compact from the previous summary'' rule, now with a reason.

\gives A one-page theorem that converts an architectural opinion (the SLM sidecar; two digest pipelines) into a mathematical necessity, upgrades Thm.~I.7.1 from an example to a characterization, and quantifies the cost of ignoring it via additive bit floors. Lands in Chapter~I as \S I.8.2's missing second half.

\subsection{A2. Deriving the regret head (and the surfacing threshold) from decision theory}
\lift{1--2 days}

\form The volume asserts $\mathrm{regret}(x) = \mathrm{stakes}(x)\times\mathrm{irreversibility}(x)\times\mathrm{anomaly}(x)$ as a design choice. It is derivable. Let $C_{\mathrm{miss}}(x) = \mathrm{stakes}(x)\cdot\mathrm{irreversibility}(x)$ be the realized loss if $x$ is bad and unattended, and let the anomaly score be a \emph{calibrated posterior} $\Pr(\mathrm{bad}\mid \text{evidence about } x)$. Then expected loss of not surfacing is exactly $C_{\mathrm{miss}}(x)\cdot\Pr(\mathrm{bad}\mid x)$ --- the multiplicative form, no longer asserted but exact, \emph{iff} anomaly is calibrated probability rather than an arbitrary score.

\proofpath Direct expected-cost minimization: surface $x$ iff $C_{\mathrm{miss}}(x)\Pr(\mathrm{bad}\mid x) \ge c_{\mathrm{att}} + C_{\mathrm{fa}}\Pr(\mathrm{good}\mid x)$, where $c_{\mathrm{att}}$ is the operator's marginal attention cost. Rearranged, this is the classical SDT likelihood-ratio criterion $\beta = \frac{1-\pi}{\pi}\cdot\frac{C_{\mathrm{fa}}}{C_{\mathrm{miss}}}$ that Def.~I.5.1 gestures at --- now derived rather than cited, with the queue's forced-zoom rule falling out as the optimal policy. Two corollaries do real work: (i) \emph{miscalibration correction} --- if the raw anomaly score $s$ is not calibrated, the correct head is $C_{\mathrm{miss}}\cdot g(s)$ with $g$ the empirical calibration map, fittable from the audit log (a measurable object, per Appendix B.3.6); (ii) the Goodhart surface sharpens --- agents attack $g$, not the formula, so the defense budget belongs on calibration integrity, unifying with A6/B2's inspection accounting.

\gives Equation~I.3 stops being a heuristic. The attention queue becomes an optimal alerting policy with a named optimality criterion, a calibration obligation, and a place where the empirical program (instrument false alarms and misses first) plugs directly into the math.

\subsection{A3. $\varepsilon$-conservation for the clean-room release ledger}
\lift{1 day, plus imports}

\form Model the release ledger as a state machine: state $=$ (cumulative spend $\sigma$, append-only log $\Lambda$); the only spending transition is $\mathrm{release}(\varepsilon_i)$, which atomically appends $(\varepsilon_i, \text{artifact hash}, \text{policy hash})$ and sets $\sigma \mathrel{+}= \varepsilon_i$; the gate refuses when $\sigma + \varepsilon_i > \varepsilon_{\max}$.

\proofpath Induction on the transition sequence gives $\sigma = \sum_{\Lambda} \varepsilon_i$ invariantly --- the identical proof shape as the bond-conservation theorem (VI.10.1), including the concurrent-invocation lemma (no over-spend under two racing releases, by the single-writer serialization the kernel already provides). Then import composition: sequential composition (budgets add) for the baseline, advanced composition ($\varepsilon' = \sqrt{2k\ln(1/\delta')}\,\varepsilon + k\varepsilon(e^\varepsilon - 1)$) when the engagement is long and the tighter accounting is worth its bookkeeping. The one honest caveat to state in-line: conservation governs \emph{recorded} spend; the claim that all spend is recorded is exactly the complete-mediation assumption (B3), which is why the hypervisor is a prerequisite, not a nicety.

\gives The clean-room's headline invariant, proved with machinery the volume already owns --- a structural rhyme (the harbor conserves one more scarce quantity through buckets) that makes the Derek/Erin product auditable rather than promissory.

\subsection{A4. Canary detection power and sequential leak testing}
\lift{1 day}

\form Plant $c$ canaries (unique sentinel spans) uniformly in Derek's corpus. The gate's verbatim/embedding suppressor has per-canary false-negative rate $\beta$. A leak carrying $k$ canary-bearing spans is detected with probability $\ge 1 - \beta^k$; solve for the canary density needed so that a leak of economic size $V$ carries $k(V)$ canaries in expectation, giving a published operating curve $\Pr(\text{detect}) = f(V)$.

\proofpath Independence across canaries under uniform planting gives the $1-\beta^k$ bound directly; the refinement worth writing is the \emph{sequential} version --- Wald's SPRT on the stream of gate outputs against the permutation null (no-leak), yielding expected time-to-detection as a function of leak rate, with the standard SPRT optimality property. This turns ``we have canaries'' into ``here is the detection-power curve and the expected latency at each leak intensity,'' the same move Chapter~VI makes for bonds.

\gives Assurance as a purchasable quantity on the leakage side: the clean-room contract can \emph{quote} detection probability, and the exfiltration bond (which cannot cover unbounded breach damage) is re-scoped to the honest job of funding detection and response --- the layer ordering the volume already believes, now with numbers.

\subsection{A5. The directory-routed message load bound}
\lift{2 days}

\form Agents generate messages at rate $\lambda$; relevance of a message to agent $i$ follows a concentration law (Zipf with exponent $\alpha > 1$ over ranked recipients is the natural model, since most coordination messages concern few parties). Broadcast delivers $\Theta(n\bar{s})$ tokens per agent per epoch. Directory-routed delivery with fan-out cap $f$ delivers $O(f\bar{s})$; the theorem quantifies the recall price: to guarantee expected relevant-message recall $\ge 1-\epsilon$, the required fan-out is $f^\star = \Theta(\epsilon^{-1/(\alpha-1)})$ --- a constant in $n$.

\proofpath Tail-mass computation on the Zipf law (the mass beyond rank $f$ is $\Theta(f^{1-\alpha})$; invert). The honest bookkeeping: the $\log n$ in my earlier note belongs to directory \emph{lookup} cost, not delivered load; separate the two lines. Add the adversarial rider from the review: since inbound messages are attacker-influenced context, the bound should be stated over \emph{provenance-labeled} tokens, and the injection defense (quoted-channel rendering) costs a constant factor, not asymptotics.

\gives Chapter II's bus gets its missing scaling claim: messaging that stays $O(1)$ in swarm size at fixed recall, with the recall/fan-out dial exposed --- the read-poverty thesis extended from discovery to communication, in one page.

\subsection{A6. Reputation-inheritance conservation for distilled skills (the no-mint theorem)}
\lift{2--3 days; the design choice is the content}

\form Let witnessed outcomes form the leaves of an evidence DAG whose internal nodes are derived artifacts (CDM-distilled skills, revisions, merges). Assign each outcome $e$ a graded value $q(e)$ and a \emph{derivation budget}: weights $w(e \to a) \ge 0$ over all artifacts $a$ derived from $e$, constrained by $\sum_a w(e\to a) \le 1$. An artifact's inherited prior is $\iota(a) = \sum_e w(e\to a)\, q(e)$, recursively through the DAG with budgets consumed at each hop.

\proofpath The theorem: total inherited prior over the DAG's closure never exceeds total witnessed value, $\sum_a \iota(a) \le \sum_e q(e)$, and no distillation cycle mints value. Proof by a potential function --- each outcome carries at most unit budget, each derivation strictly consumes it, so the sum is a supermartingale under any sequence of distillations; cycles are safe because re-derivation from an artifact draws on the artifact's (already-discounted) budget, not the leaf's. The crux is the design decision the proof forces into the open: budgets are \emph{per outcome, split across all derivations}, not per derivation --- otherwise the distill-loop mints arbitrarily (derive $k$ skills from one episode, each inheriting full value).

\gives Closes Gap~\#4 (skill versioning cold-start) safely: v2 warm-starts from provenance, fraud cannot amplify through distillation chains, and the CDM pipeline (Mechanism~6) gets its economic safety proof. Also the estate's export/import ceremony inherits the same conservation for free.

\subsection{A7. The Theorem I.6.1 falsification experiment (executable)}
\lift{hours to a first figure; 1--2 days to a clean writeup --- I can run this in the container on request}

\form Simulate $N$ artifacts with a hidden critical set of size $k$; construct digests at token budgets $B$ sweeping through the information floor $\log_2\binom{N}{k} - \log_2\binom{m}{k}$: (i) an optimal covering-design digest at small $N$, (ii) greedy set-cover digests, (iii) a learned scorer over noisy features. Measure empirical miss rate versus $B$ and overlay the floor.

\proofpath Not a proof --- a falsification apparatus. The prediction is sharp: zero-miss is unattainable below the floor for \emph{any} digest scheme, and the empirical curves should hug it from above, with the gap between greedy/learned and the floor measuring the price of realistic (non-oracle) digest construction. The second panel is the adversarial variant from Solution 6.5$^\star$: corrupt a fraction $\phi$ of features and watch the effective floor inflate --- the first empirical trace of the mimicry/inflation attack. The third panel closes the loop with A1: measure the joint-digest penalty when two readers' critical sets are disjoint versus overlapping.

\gives The volume's first datum, its best figure, and a public artifact (the beginning of the coordination benchmark the review argued you should own). It also stress-tests the theorem's modeling: if real digests systematically beat the floor, the critical formalization is wrong, and finding that out cheaply is the point of having stated a falsifiable bound.

\section{Tier B --- roughly a week each}

\subsection{B1. The digest--zoom Pareto frontier: two-stage group testing and one-sided rate--distortion}
\lift{1 week; the derivations are standard in technique, the assembly is new}

\form The digest writer knows the state $X^N$ ($X_i \in \{$critical, inert$\}$, i.i.d.\ Bernoulli$(p)$ in the average-case model); the operator does not. A digest is $B$ bits about $X^N$; the operator opens the flagged set and may then \emph{zoom} adaptively. Total operator cost: $c_{\mathrm{read}}B + c_{\mathrm{open}}\mathbb{E}|\hat S| + c_{\mathrm{miss}}\mathbb{E}|S\setminus\hat S|$. The right information-theoretic object is a \emph{two-constraint} rate--distortion program:
\[ R(\delta, f) \;=\; \min_{P_{\hat X\mid X}} I(X;\hat X) \quad \text{s.t.}\quad \Pr(X{=}1, \hat X{=}0) \le \delta,\;\; \Pr(\hat X{=}1) \le f, \]
false negatives priced by $\delta$ (misses), false positives priced only through the open budget $f$. Your Thm.~I.6.1 is the $(\delta{=}0, f{=}m/N)$ worst-case corner of this program.

\proofpath Three deliverables. (i) Solve $R(\delta,f)$ in closed form for the binary source via the Lagrangian/test-channel method --- a finite convex program whose KKT conditions I can carry by hand; verify the corners ($f{=}1 \Rightarrow R{=}0$: flag everything; $\delta{=}0, f{=}p \Rightarrow R \to h(p)$-scale: the digest must localize). (ii) The \emph{zoom advantage theorem}: adaptive second-stage inspection with group-level summaries (zooming a directory before its files) is Hwang-style generalized binary splitting inside the flagged region, reducing expected opens from $fN$ to $k\log_2(fN/k) + O(k)$ whenever group verification is available --- quantifying, for the first time, what the ``zoom'' in digest-with-zoom is mathematically worth over a flat digest. (iii) Assemble (i)+(ii) into the Pareto frontier (digest bits) $\times$ (expected opens) at fixed miss tolerance, with the operator's optimum at the price-ratio tangent point $c_{\mathrm{read}}/c_{\mathrm{open}}$.

\gives Chapter I's \S I.6 grows from one worst-case bound into a small theory: average-case floors, the formal value of adaptivity, and a knob (the tangent point) that turns operator time-accounting into digest sizing. It also gives A7's experiment its second axis, and it is the strongest single candidate for a standalone paper --- the mathematical identity of ``summarize the state, link to the work.''

\subsection{B2. The inspection tower: critical audit rates, contraction, and reputation as amortized verification}
\lift{1--1.5 weeks}

\form Level $k$ is an inspection game: the judge cheats with probability $q$ for one-shot gain $G_k$; the inspector audits with probability $\rho_k$ at cost $a$ per audit, detecting with probability $d_k$ and slashing bond $B_k$. Levels stack: level $k{+}1$ audits level-$k$ auditors, sampled sealed from a pool spanning $\ge q_{\mathrm{cl}}$ disjoint principal-cliques. Above the tower sits a reputation model in the Kreps--Wilson/Fudenberg--Levine mold: a fraction $\mu_0$ of judges are commitment (honest) types; strategic types decide whether to mimic.

\proofpath Four steps, each standard in technique. (i) The stage game's committed-inspection solution: deterrence at $\rho_k^\star = G_k/(d_k B_k)$, minimal audit spend $a\,G_k/(d_k B_k)$ per judged outcome --- Becker's condition, now with the audit \emph{cost} accounted, which the conjecture omits. (ii) Contraction: under sealed sampling from $q_{\mathrm{cl}}$ cliques, a briber corrupting level-$k{+}1$ must in expectation pay off order $q_{\mathrm{cl}}$ independent auditors' expected forfeitures before capturing the audit, so $G_{k+1} \le G_k\lambda$ with $\lambda = O(1/q_{\mathrm{cl}})$; then required bonds $B_k \gtrsim G_0\lambda^k/(\rho d)$ sum geometrically --- \emph{finite total bond capitalization certifies an infinite tower}, which is the contraction Conj.~III.11.1 asks for, with its two assumptions (non-shrinking stakes, sealed multi-clique sampling) named as the load-bearers. (iii) Termination alternative: a finite tower rooted at an exogenously honest oracle (the operator at $n{=}1$) with the same inequality at each level. (iv) The reputation layer: derive the declining audit schedule $\rho_t$ that keeps the strategic type's IC binding as the posterior $\mu_t \to 1$ --- audits fall as continuation value rises --- and prove lifetime audit spend per honest judge is $O(1)$ versus $\Theta(T)$ under flat auditing. That last result \emph{is} ``reputation is amortized verification,'' stated as a theorem.

\gives The volume's central open conjecture becomes a parameterized theorem plus a measurement obligation (does the deployed pool satisfy sealed multi-clique sampling?), the audit budget becomes a designed quantity rather than a hope, and Chapter III gains its deepest one-liner with a proof behind it.

\subsection{B3. Hypervisor enforceability is Ramadge--Wonham controllability}
\lift{1 week}

\form Model agent behavior as a prefix-closed language $L$ over an alphabet $\Sigma = \Sigma_c \cup \Sigma_u$: \emph{controllable} events $\Sigma_c$ are those the hypervisor mediates pre-commit (fs writes, network egress, exec), \emph{uncontrollable} events $\Sigma_u$ are everything it cannot intercept (model-internal steps, in-context content). A policy is a safety specification $K \subseteq L$.

\proofpath Direct adaptation of supervisory control theory: a supervisor achieving exactly $K$ exists iff $K$ is controllable --- $\overline{K}\Sigma_u \cap \overline{L} \subseteq \overline{K}$, i.e., no uncontrollable event escapes the specification's closure. The theorem for the volume: \textbf{Port Daddy's hypervisor can \emph{regiment} policy $P$ iff $P$ is a prefix-closed safety language controllable with respect to $\Sigma_u$}; everything else is detect-and-compensate, by your own Thm.~II.8.1. The proof is the Ramadge--Wonham construction specialized to the daemon's mediation map, plus one wrinkle I will treat honestly: policies referencing unmediated \emph{state} (not just events) need the partial-observation extension (observability/normality conditions), which strictly shrinks the regimentable set and should be stated as such. The corollary is a table the ledger has been missing: force-push (network event) --- regimentable; out-of-scope write, deleting a failing test (fs events over policy-expressible paths) --- regimentable; ``confident lie in a report'' (semantic content, uncontrollable) --- forever detect-and-compensate. Re-grade every affected ledger row against the alphabet.

\gives The hypervisor stops being an engineering preference and becomes a theorem about which promises the product may make: an exact, citable boundary between \emph{regimented} and \emph{enforced}, resolving \S I.4.4's open problem (footgun guards as structural authority) for every policy inside the controllable set --- and honestly fencing what stays outside it.

\subsection{B4. A tractable deontic-conflict fragment for the condominium harbor}
\lift{1 week}

\form Commitment language $\mathcal{L}_c$: (a) ground Horn clauses over a finite atom vocabulary (invariants, accepted obligations); (b) interval resource claims $(\mathrm{res}, [t_1,t_2], \mathrm{owner}, \mathrm{excl})$; (c) difference constraints linking commitment deadlines. A \emph{conflict} is a deontic inconsistency: derivable $\bot$, overlapping exclusive intervals, or an infeasible constraint cycle.

\proofpath Three results. (i) Membership: consistency in $\mathcal{L}_c$ is decidable in polynomial time and \emph{incrementally} --- Horn by watched-literal unit propagation (linear in occurrence size per update), intervals by interval trees ($O(\log n)$ per insert), difference constraints by negative-cycle detection on the constraint graph; I will write the combined incremental algorithm and its amortized bound. (ii) The frontier: one step outside --- arbitrary clauses, or disjunctive obligations $O(a \vee b)$ --- makes consistency NP-complete (reduction from SAT, immediate), so the language design is not taste but the exact price of staying detectable-at-proposal-time. (iii) Completeness-with-regimentation: within $\mathcal{L}_c$, every conflict is detected \emph{before} acceptance, and since acceptance is a mediated write, B3 makes conflicting acceptances \emph{regimentable} --- rejected pre-commit, extending Thm.~VI.5.1's advisory completeness to structural prevention in the shared-harbor setting.

\gives The condominium chapter's algorithmic core: the lookout's contradiction-scan becomes a polynomial incremental checker with a proved completeness property, ``who decides what conflicts are'' splits cleanly into mechanical detection and chartered resolution, and the schema designer gets a theorem telling them exactly which expressive conveniences would forfeit tractability.

\subsection{B5. Engine substitution: Akerlof unraveling inside one identity, and resurrection soundness}
\lift{4--6 days}

\form Two engine qualities $\theta_H > \theta_L$ with costs $c_H > c_L$; an identity carries reputation $r$ earned under some engine history; buyers observe $r$ but, absent attestation, not the current engine. A strategic principal may substitute $\theta_L$ after building $r$ on $\theta_H$.

\proofpath Two-period model, standard lemons technique. Without attestation, buyer beliefs price the identity at $\mathbb{E}[\theta \mid r, \text{equilibrium strategy}]$; substitution is profitable at any pooled price above $\theta_L$'s cost, so pooling on $\theta_H$ cannot survive: either prices discount to the mixture (taxing honest $\theta_H$ sellers, who exit or substitute --- the ratchet to $\theta_L$) or the market unravels --- Akerlof's argument run \emph{within} a single identity, which is the novel framing. With daemon-attested engine identifiers on every witnessed outcome, reputation conditions on $(\text{person}, \text{engine class})$, the information asymmetry closes, and I then prove \emph{resurrection soundness}: under (i) credential and continuity-witness verification, (ii) engine attestation recorded in successor outcomes, and (iii) open commitments closed or renegotiated with escrow delta, the sanction-respecting property (Def.~III.6.1) is preserved across provider migration --- the migration cannot be used to shed either sanctions or engine history. The same theorem, read sideways, closes skill-versioning Gap~\#4: reputation attaches to (identity, substrate) pairs with disclosed substrate, one rule for both holes.

\gives The resume-button feature ships with its own security theorem and its own named attack; Chapter III gains the market-microstructure result it currently lacks; and the buyer-facing story becomes crisp: reputation you can trust \emph{because} the engine behind it is attested.

\subsection{B6. The newcomer ramp as screening: the front-loaded probation theorem}
\lift{3--5 days; the result is a short exchange argument with a surprising sign}

\form Two types enter as newcomers: honest (discount $\delta_h$) and fraud-planning (effective discount $\delta_f < \delta_h$, since their horizon truncates at the planned reset). The designer chooses the ceiling schedule $\{\kappa_t\}$, equivalently the per-period surplus gaps $g_t = \pi(\kappa^\star) - \pi(\kappa_t) \ge 0$. Deterrence requires the \emph{fraudster's} discounted ramp cost to exceed the maximal one-shot gain: $\sum_t \delta_f^t g_t \ge G_{\max}$. The design objective: minimize the \emph{honest} newcomer's discounted burden $\sum_t \delta_h^t g_t$ subject to that constraint.

\proofpath The ratio of deterrence-value to honest-cost of a unit gap at time $t$ is $(\delta_f/\delta_h)^t$, strictly decreasing in $t$. An exchange argument --- move gap mass from any later period to any earlier one, holding the deterrence constraint tight --- strictly lowers honest cost; iterate to the corner. Conclusion: the optimal schedule is \emph{maximally front-loaded}: the deepest ceiling restriction immediately, released as fast as the deterrence constraint permits --- a probation cliff, not a long slow ramp. The result is mildly counter-intuitive (deferred-compensation instincts, \`a la Lazear, suggest back-loading) and the resolution is instructive to write out: here the punishment \emph{is} the ramp re-paid after reset, so it must land where the short-horizon type still feels it --- early. Extensions I can carry in the same draft: continuum of types (the cliff persists), and the comparative static that as fraud gains $G_{\max}$ grow, the cliff deepens before it lengthens.

\gives Theorem III.6.2 currently says ``choose the schedule so $W$ exceeds the fraud gain''; this pins down \emph{which} schedule, with a clean behavioral prediction (probation cliffs dominate gradual ramps) that a deployment can implement tomorrow and A/B against churn.

\subsection{B7. The costly-escalation signaling equilibrium and the tuning band for the debit}
\lift{4--6 days}

\form Agents privately observe urgency $u \sim F$; escalating surfaces the item; the operator validates iff genuinely urgent (revealed on attention). A dismissed escalation debits standing by $\delta$, worth $w$ per unit in future surfacing weight; escalating when urgent yields benefit $b(u)$, increasing.

\proofpath Single-crossing: expected payoff of escalating, $b(u) - \delta w\,(1 - V(u))$ with validation probability $V(u)$ increasing, crosses zero once, so the equilibrium is a threshold $u^\star(\delta)$ solving $b(u^\star) = \delta w (1 - V(u^\star))$ --- a separating equilibrium in the Spence mold, with the queue's ``crying wolf is costly'' realized as money-burning that only high types rationally risk. Comparative statics give the tuning band: $u^\star$ increases in $\delta$; alarm load is $1 - F(u^\star)$; false silence (genuine distress suppressed) is $F(u_{\mathrm{crit}}) - F(u^\star)$ truncated at zero. The deliverable is the two-sided constraint $\delta \in [\delta_{\min}, \delta_{\max}]$: below $\delta_{\min}$ the operator's alarm-fatigue budget (from A2's attention accounting) is violated; above $\delta_{\max}$ the miss-risk on the distress lane exceeds the SDT criterion. Both endpoints are computable from instrumented quantities the roadmap already plans to collect.

\gives Section I.7.3's mechanism gets an equilibrium existence proof, an explicit operating band for its one free parameter, and a bridge to the empirical program: $\delta$ is fit from measured validation rates, not chosen by vibes.

\subsection{B8. The specialization boundary for sole-responsibility roles}
\lift{3--4 days}

\form Function-$F$ requests arrive Poisson($\lambda_F$). Option S: a sole specialist, M/M/1 with service rate $\mu_s$. Option P: fold the traffic into a pool of $c$ generalists, M/M/$c$ at rate $\mu_g \le \mu_s$ each (the specialist's rate advantage is the skill premium). Sole ownership additionally yields accountability value $A$ per period (clean attribution, per-role ledger, succession clarity).

\proofpath Exact comparison via Erlang-C: closed form for $c=2$, the general condition as an inequality in the Erlang-C delay probability. The theorem is the boundary: sole ownership dominates iff $\mu_s/\mu_g \ge g(\rho, c)$ or the latency deficit is smaller than $A$ priced at the delay cost --- with $g$ explicit, and the known queueing folklore (pooling wins on variance) now carrying its exceptions in the open: high skill premia, low utilization, or high attribution value. A worked corner I will include because it matters operationally: the specialist as single point of failure --- M/M/1 with breakdowns (the death/heartbeat model the volume already has) --- which quantifies the succession-rule requirement rather than leaving it as hygiene.

\gives Mechanism 2 (calendar avatar, test writer, roadmap owner) gets an honest boundary instead of an aesthetic: when the pattern pays, when it silently costs latency, and what the succession rule is worth --- all in quantities the daemon already meters.

\subsection{B9. Context paging with an untrusted pin oracle}
\lift{4--5 days; mostly a disciplined import with one new term}

\form Resident context is a cache of size $k$ tokens over spans in the durable buffer; span requests follow the agent's action sequence; eviction is online. Critical spans are \emph{pinned} by a feature oracle that an adversary corrupts at rate $\phi$ (mis-pin or missed pin).

\proofpath Import the weighted-caching result (Landlord): $k/(k-h+1)$-competitive against offline with cache $h$, applied to the unpinned region --- variable span sizes are exactly why the weighted/file-caching version is needed rather than vanilla LRU's Sleator--Tarjan bound, and I will keep that distinction honest. The new content is the degradation term: a charging argument that mis-pins cost at most their carried size (wasted residency) and missed pins cost at most one re-fetch each, giving total cost $\le$ Landlord bound on the honest region $+\ \phi N \cdot c_{\mathrm{refetch}}$ --- linear, graceful, and tied by the same $\phi$ to Solution 6.5$^\star$'s forgeable-feature budget, so the paging layer and the digest layer share one adversarial parameter.

\gives Mechanism 3's buffer stops being plumbing: a competitive guarantee for ``what stays in context,'' a single named adversarial knob shared with the legibility theory, and the formal restatement of the effective-context trap --- paying for tokens that competitive analysis says should have been evicted.

\section{Tier C --- multi-week}

\subsection{C1. Conditional noninterference for the clean room via unwinding}
\lift{2--3 weeks; standard technique, but the system model must be built with care}

\form Domains $\{D, E, G_{\to E}, G_{\to D}, \mathrm{pub}\}$ with an \emph{intransitive} flow policy: $D \rightsquigarrow G_{\to E} \rightsquigarrow E$ and $E \rightsquigarrow G_{\to D} \rightsquigarrow D$, with the direct edges $D \rightsquigarrow E$ and $E \rightsquigarrow D$ absent. The system is the hypervisor-mediated transition machine: agent steps, gate steps, ledger steps.

\proofpath Rushby's intransitive-noninterference framework: define the purge function for the policy, state the security condition (observations of $E$ depend only on sources permitted via the gate chain), and discharge it by the three unwinding conditions --- output consistency, weak step consistency, and local respect --- proved per transition class. Complete mediation (B3) discharges the ``no other edges'' premise; gate correctness is assumed, and I will state it as the residual oracle it is (the clean room \emph{relocates} trust into two auditable declassifiers; it does not eliminate it). The honest scope paragraph writes itself and must be written: the guarantee is possibilistic --- timing, token-count, and termination channels are outside the model, as is a human reading gate output and remembering it; the theorem bounds \emph{channels}, not minds.

\gives The Derek/Erin design's central promise as a theorem with named unwinding obligations --- each one a checkable property of a concrete daemon transition --- which is precisely the form in which a security claim can later be mechanized (the same ProVerif-adjacent discipline Chapter V already practices, extended from secrecy of keys to flow of labeled data).

\subsection{C2. Buying your way out of the Myerson--Satterthwaite corner: the minimum viable take-rate}
\lift{2 weeks in a stylized type model}

\form The harbor is not a bilateral mechanism: it has an insurer/platform side with fee income $\Phi$ (bond carry, listing, settlement fees). With an outside budget, ex-post efficient trade is implementable (subsidize the VCG deficit); the question MS actually leaves open here is \emph{how much} budget.

\proofpath In the canonical uniform-types model, compute the expected VCG deficit of efficient bilateral trade in closed form (a standard exercise I can carry, and one worth re-deriving cleanly rather than citing), then state the theorem: the platform escapes the MS corner iff expected fee income per trade weakly exceeds expected deficit, giving a \emph{minimum viable take-rate} $\tau^\star$ as an explicit function of the value distributions. Corollaries: (i) the take-rate is the price of ex-post efficiency, a sentence Chapter IV can put on a slide; (ii) with the common-value contamination the review flagged (shared grading oracle correlates values), the relevant distortion shifts from MS to lemons, and the same fee budget is better spent on information design --- disclosure of reputation vectors --- than on subsidies, which I will state as a scoped comparison rather than a grand unification.

\gives Chapter IV's most cited impossibility stops being a corner the design nervously respects and becomes a budget line the platform prices --- and the analysis says when subsidy is the right instrument and when disclosure is.

\section{What I decline to claim}
In the ledger's spirit: \textbf{gossip convergence under real partitions} is an empirical distribution, not a theorem worth having beyond the stated model (the volume already grades this correctly). \textbf{Cross-operator attestation} is a deployment-and-hardware problem; I can specify the TEE/key-transparency evaluation, not prove the gap closed. \textbf{Folk-theorem cartel bounds at the federation layer} are provable only in monitoring models stylized past usefulness; I would write the model-selection memo, not claim the bound. \textbf{The categorical/sheaf program} beyond the toy gluing model (already solved in the prior document) adds vocabulary, not power, at current scale --- I stand by commit-or-cut. And every result above that touches deployment parameters ($d$, $\beta$, $\phi$, calibration maps, validation rates) is a theorem \emph{conditional on measurement}; the instrumentation plan of Appendix B.3.6 is the co-requisite, not an afterthought.

\section{Suggested batching}
Week 1: A1 $+$ A2 $+$ A7 (the theorem, its decision-theoretic twin, and the experiment --- one coherent legibility package). Week 2: B3 (hypervisor characterization; unblocks the most ledger rows and C1). Week 3: B2 (inspection tower). Week 4: B1 (the frontier paper). Then B5$+$B6 as the Chapter III pair, B4 for the condominium, and C1/C2 as the long arcs. A1$+$A2$+$A7$+$B1 together are a self-contained paper: \emph{``The Price of a Summary: Information Floors, Adaptive Zoom, and Split Heads for Operator Legibility.''}

\end{document}
